% !TeX root = ../../Book5.tex
% 纲要标题：## 附录 C　实、热带与组合代数几何
% 纲要位置：几何学计划纲要.md，第 6375 行。
% 正文状态：未撰写；本文件目前是框架。

% 纲要内容开始
% 
% > 以实不等式和赋值退化为两条独立入口，保持各自的对象、拓扑和重数。
% 
% * **实点与半代数结构**
%   * 实点集和 Zariski 闭包
%   * Tarski–Seidenberg 的准确陈述
%   * 实曲线的连通分支
%   * 正性证书和复几何的区别
% 
% * **初始退化与热带化**
%   * 赋值、权和初始理想
%   * Gröbner 退化与平坦族
%   * 热带超曲面的计算
%   * 平衡条件与重数
%   * 平凡及非平凡赋值的区别
% 
% * **组合和枚举**
%   * 格多面体与次数
%   * 热带曲线的有限模型
%   * 对应定理的具体假设
%   * 组合计数不能无条件替代代数重数
% 
% ---
% 
% 纲要内容结束

\chapter{实、热带与组合代数几何}
\label{g5:app:C}
\BookChapterNumber{g5:app:C}

待撰写

\section{实点与半代数结构}
\label{g5:sec:C01}
待撰写

\subsection{实点集和 Zariski 闭包}
\label{g5:subsec:C01:01}
待撰写

\subsection{Tarski–Seidenberg 的准确陈述}
\label{g5:subsec:C01:02}
待撰写

\subsection{实曲线的连通分支}
\label{g5:subsec:C01:03}
待撰写

\subsection{正性证书和复几何的区别}
\label{g5:subsec:C01:04}
待撰写

\section{初始退化与热带化}
\label{g5:sec:C02}
待撰写

\subsection{赋值、权和初始理想}
\label{g5:subsec:C02:01}
待撰写

\subsection{Gröbner 退化与平坦族}
\label{g5:subsec:C02:02}
待撰写

\subsection{热带超曲面的计算}
\label{g5:subsec:C02:03}
待撰写

\subsection{平衡条件与重数}
\label{g5:subsec:C02:04}
待撰写

\subsection{平凡及非平凡赋值的区别}
\label{g5:subsec:C02:05}
待撰写

\section{组合和枚举}
\label{g5:sec:C03}
待撰写

\subsection{格多面体与次数}
\label{g5:subsec:C03:01}
待撰写

\subsection{热带曲线的有限模型}
\label{g5:subsec:C03:02}
待撰写

\subsection{对应定理的具体假设}
\label{g5:subsec:C03:03}
待撰写

\subsection{组合计数不能无条件替代代数重数}
\label{g5:subsec:C03:04}
待撰写
